%(BEGIN_QUESTION)
% Copyright 2011, Tony R. Kuphaldt, released under the Creative Commons Attribution License (v 1.0)
% This means you may do almost anything with this work of mine, so long as you give me proper credit

En digital nivåtransmitter har et kalibrert inngangsområde fra 0 til 40 tommer, og en 10-bits utgang (område 0 til 1023 ``count''). Fullfør følgende verditabell for denne transmitteren, med antatt perfekt kalibrering (ingen feil):

% No blank lines allowed between lines of an \halign structure!
% I use comments (%) instead, so that TeX doesn't choke.

$$\vbox{\offinterlineskip
\halign{\strut
\vrule \quad\hfil # \ \hfil & 
\vrule \quad\hfil # \ \hfil & 
\vrule \quad\hfil # \ \hfil & 
\vrule \quad\hfil # \ \hfil \vrule \cr
\noalign{\hrule}
%
% First row
Inngangsnivå & Prosent av spenn & Count & Count \cr
%
% Another row
(tommer) & (\%) & (desimal) & (heksadesimal) \cr
%
\noalign{\hrule}
%
% Another row
 & 37 &  &  \cr
%
\noalign{\hrule}
%
% Another row
 & 82 &  &  \cr
%
\noalign{\hrule}
} % End of \halign 
}$$ % End of \vbox

\underbar{file i01621}
%(END_QUESTION)





%(BEGIN_ANSWER)

Full uttelling gis for å ha én av de alternative verdiene i hver celle i tabellen (studenten trenger altså ikke angi {\it begge} count-verdiene som vises i hver celle):

% No blank lines allowed between lines of an \halign structure!
% I use comments (%) instead, so that TeX doesn't choke.

$$\vbox{\offinterlineskip
\halign{\strut
\vrule \quad\hfil # \ \hfil & 
\vrule \quad\hfil # \ \hfil & 
\vrule \quad\hfil # \ \hfil & 
\vrule \quad\hfil # \ \hfil \vrule \cr
\noalign{\hrule}
%
% First row
Inngangsnivå & Prosent av spenn & Count & Count \cr
%
% Another row
(tommer) & (\%) & (desimal) & (heksadesimal) \cr
%
\noalign{\hrule}
%
% Another row
14.8 & 37 & 378 eller 379 & 17A eller 17B \cr
%
\noalign{\hrule}
%
% Another row
32.8 & 82 & 838 eller 839 & 346 eller 347 \cr
%
\noalign{\hrule}
} % End of \halign 
}$$ % End of \vbox

%(END_ANSWER)





%(BEGIN_NOTES)

{\bf Dette spørsmålet er kun ment for eksamen, ikke arbeidsark!}.

%(END_NOTES)
